\documentclass[a4paper,11pt]{article}
\usepackage[french]{babel}
\usepackage{amsmath}
\usepackage[osf,sups]{Baskervaldx} % lining figures
\usepackage[bigdelims,cmintegrals,vvarbb,frenchmath,baskervaldx]{newtxmath} % math font
\usepackage{graphicx}
\usepackage{xcolor}
\usepackage{url}
\usepackage[hidelinks]{hyperref}

% hyperrref configuration: remove ugly boxes and colors, keep urls in blue
\hypersetup{colorlinks,linkcolor=black,citecolor=black,urlcolor={blue!80!black}}


% paragraphs
\setlength{\parindent}{0em}
\setlength{\parskip}{1em}

% margins
\addtolength{\hoffset}{-3em}
\addtolength{\textwidth}{6em}
\addtolength{\voffset}{-3em}
\addtolength{\textheight}{6em}

% macros
\def\ud{\mathrm{d}}
\def\R{\mathbf{R}}
\def\K{\mathbf{K}}
\def\S{\mathcal{S}}
\def\Z{\mathbf{Z}}

\begin{document}

\begin{center}\LARGE
stage L3 : plan de travail semaines 1,2,3
\end{center}

{\bf Notations.}
L'\underline{espace de Schwartz}~$\S(\R^d)$ ou simplement~$\S$ est l'ensemble
des fonctions~$\mathcal{C}^\infty$ telles que toutes leurs dérivées sont à
décroissance rapide
\[
	f\in\S(\R)
	\iff
	\forall p,n : \lim_{|x|\to\infty}\frac{\ud^n}{\ud x^n}f(x)=0
\]
et la même condition sur les dérivées partielles pour~$f\in\S(\R^d)$.

On utilise la convention suivante pour la \underline{transformée de Fourier}
en~$\R^d$:
\[
	\widehat f(y)=\int_{\R^d}e^{-2\pi i x\cdot y}\ud x
\]
définie d'abord dans~$\S$ ou~$L^1$ puis éténdue à~$L^2$ par continuité.

Les \underline{distributions tempérées} sont le dual topologique de~$\S$.
Ainsi, une distribution tempérée est un opérateur linéaire continu~$T:\S\to\K$,
noté~$(T,\varphi)$ ou parfois~$\int T\varphi$.  Cette dernière notation est
justifiée du fait que~$\S'$ contient toutes les fonctions localement
intégrables à croissance polyomiale, via l'idéntification
injective~$\left(f,\varphi\right)=\int f(x)\varphi(x)\ud x$.

Un~\underline{réseau} de~$\R^2$ est un subgroupe additif généré par deux
vecteurs linéairement indépendants.  L'exemple canonique est~$\Z^2$, généré par
les vecteurs $\left(\begin{smallmatrix}1\\0\end{smallmatrix}\right)$
et~$\left(\begin{smallmatrix}0\\1\end{smallmatrix}\right)$.  Notez que
un même réseau peut être engendré par bases différentes, par
exemple~$\Z^2$ est aussi générée par les vecteurs
$\left(\begin{smallmatrix}1\\0\end{smallmatrix}\right)$
et~$\left(\begin{smallmatrix}10\\1\end{smallmatrix}\right)$.
On verra plus tard que toutes les bases d'un même réseau ont le même
determinant.


{\bf Exercice 0. (transformée au sens des distributions)}
La transformée de Fourier d'une distribution tempérée~$T\in\S'$ est la
distribution tempérée~$\widehat T\in\S'$ définie par~$\left(\widehat
T,\varphi\right)=\left(T,\widehat\varphi\right)$.  Vérifiez que~$1$, $\delta$,
$x\mapsto x^2$
et~$x\mapsto1/\sqrt{|x|}$ sont des distributions tempérées, et calculez leurs
transformées de Fourier.

{\bf Exercice 1. (formule de Poisson)}
Soit~$\varphi\in\S$, démontrez la formule de Poisson
\[
	\sum_{n\in\Z}\varphi(n)=\sum_{k\in\Z}\widehat\varphi(k)
\]
Généralisez cette formule au cas de points à distance
arbitraire~$\displaystyle\sum_{p\in\lambda\Z}\varphi(p)=\cdots$ ainsi qu'à
dimensions
supérieures~$\displaystyle\sum_{\mathbf{n}\in\Z^d}\varphi(\mathbf{n})=\cdots$.

{\bf Exercice 2. (réseau dual)}
Soit~$\Lambda$ un réseau de~$\R^2$.
Démontrez une formule de Poisson généralisée
\[
	\sum_{p\in\Lambda}\varphi(p)
	=
	c_\Lambda\sum_{q\in\Lambda^*}\widehat\varphi(q)
\]
avec une constante~$c_\Lambda$ et un réseau dual~$\Lambda^*$ que l'on
determinera explicitement à partir de~$\Lambda$.  Donnez une interprétation
géométrique pour~$c_\Lambda$ et~$\Lambda^*$.\\
{\sl Indication: appliquez un changement de variable linéaire à la formule de
Poisson classique sur~$\Z^2$.}

{\bf Exercice 3. (fomule de Poisson au sens des distributions)}
...

{\bf Exercice 4. (transformée d'une projection)}
On définit la~\emph{projection} sur l'axe~$x$ d'une fonction~$f(x,y)$ comme la
fonction~$g(y)=\int_{-\infty}^\infty f(x,y)\ud y$.  Exprimez la transformée de
Fourier de~$g$ en termes de celle de~$f$.  Généralisez cette formule à la
projectoin orthogonale sur une droite quelconque (engendrée par un vecteur
$\left(\begin{smallmatrix}\cos\alpha\\\sin\alpha\end{smallmatrix}\right)$), et à
la projection orthogonale d'une fonction de~$\R^d$ sur un sous-espace
de dimension~$k<d$.

{\bf Exercice 5. (transformée d'une coupe)}
La~\emph{coupe} horizontalle d'épaisseur~$2a$ de la fonction~$f(x,y)$ est la
fonction~$f(x,y)\mathbf{1}_{[-a,a]}(y)$.  Exprimez la transformée de la coupe
en termes de la transformée de~$f$.  Généralisez cette formule aux coupes
de~$\R^d$ autour d'un sous-espace arbitraire de dimension~$k<d$.

{\bf Exercice 6. (projection et coupe de distributions)}
Étudiez la généralisation des exercices 3 et 4 quand~$f$ n'est pas une fonction
mais une distribution ``discrète'', de la forme~$\sum_ia_i\delta(x-p_i)$

{\bf Exercice 7. (pavage apériodique)}
Considérez le peigne de Dirac bidimensionel~$P(x)=\sum_{p\in\Z^2}\delta(x-p)$.
Calculez la projection+coupe de~$P$ sur une droite d'angle~$\alpha$.  Démontrez
que la distribution résultante sur~$\R$ est périodique si et seulement
si~$\tan\alpha\in\mathbf{Q}$, et calculez son période.

\end{document}
